% Sprachüberprüfung
\RequirePackage[english, ngerman]{babel}
\ifstr{\lxlangcode}{EN}{
	\makeatletter
	\main@language{english}
	\makeatother
	%%%%%%%%%%%%%%%%%%%%%%%%%%%%%%%%%%%%%%%%%%%%%%%
%standardphrasen und schnipsel in deutsch     %
%dient als vorlage für alle anderen sprachen  %
%%%%%%%%%%%%%%%%%%%%%%%%%%%%%%%%%%%%%%%%%%%%%%%


\newcommand{\anrede} {Ladies and Gentlemen}
\newcommand{\anredefrau} {Dear Ms}
\newcommand{\anredeherr} {Dear Mr}


\newcommand{\nr} {No.}
\newcommand{\datum} {Date}
\newcommand{\kundennummer} {Customer-No.}
\newcommand{\ansprechpartner} {contact person}
\newcommand{\bearbeiter} {Bearbeiter}
\newcommand{\gruesse} {Sincerely yours, }
\newcommand{\vom} {from}
\newcommand{\von} {from}
\newcommand{\seite} {page}
\newcommand{\uebertrag} {amount carried over}


\newcommand{\position} {Pos.}
\newcommand{\artikelnummer} {Part No.}
\newcommand{\bild} {Picture}
\newcommand{\keinbild} {n/a}
\newcommand{\menge} {Q'ty}
\newcommand{\bezeichnung} {Description}
\newcommand{\seriennummer}{Serial No.}
\newcommand{\ean}{EAN}
\newcommand{\projektnummer}{Procekt No.}
\newcommand{\charge}{Charge}
\newcommand{\mhd}{Best before}
\newcommand{\einzelpreis} {Price}
\newcommand{\gesamtpreis} {Amount}
\newcommand{\nettobetrag} {net amount}
\newcommand{\schlussbetrag} {Total}

\newcommand{\weiteraufnaechsterseite} {to be continued on next page  ..}

\newcommand{\zahlung} {Payment terms:}


% angebot (sales_quotion)
\newcommand{\angebot} {Quotation}
\newcommand{\angebotsformel} {we are please to make following offer:}
\newcommand{\angebotdanke} {We thank you for your request and looking forward to receiving your order.}
\newcommand{\angebotgueltig} {This offer is valid until}		%Danach wird das Datum eingefügt, falls das grammatisch nicht funktionieren sollte müssen wir eine ausnahme für die sprache definieren
\newcommand{\angebotfragen} {If you have any questions do not hesitate to conatct us.}
\newcommand{\angebotagb} {Our general terms and conditions (AGB) apply. We will send them to you on request.}


% auftragbestätigung (sales_order)
\newcommand{\auftragsbestaetigung} {Order}
\newcommand{\auftragsnummer} {Order No.}
\newcommand{\ihreBestellnummer} {Your reference no.}
\newcommand{\auftragsformel} {hiermit bestätigen wir Ihnen folgende Bestellpostionen:}
\newcommand{\lieferungErfolgtAm} {Die Lieferung erfolgt am} %Danach wird das Datum eingefügt, falls das grammatisch nicht funktionieren sollte müssen wir eine ausnahme für die sprache definieren
\newcommand{\auftragpruefen} {Bitte kontrollieren Sie alle Positionen auf Übereinstimmung mit Ihrer Bestellung! Abweichungen teilen Sie innerhalb von 3 Tagen mit!}

% lieferschein (sales_delivery_order)
\newcommand{\lieferschein} {Lieferschein}

% rechnung (invoice)
\newcommand{\rechnung} {Rechnung}
\newcommand{\rechnungsdatum} {Rechnungsdatum}
\newcommand{\ihrebestellung} {Ihr Bestellung}
\newcommand{\lieferdatum} {Lieferdatum}
\newcommand{\rechnungsformel} {für unsere Leistungen erlauben wir uns, folgende Positionen in Rechnung zu stellen:}
\newcommand{\zwischensumme} {Zwischensumme}
\newcommand{\leistungsdatumGleichRechnungsdatum} {Das Leistungsdatum entspricht, soweit nicht anders angegeben, dem Rechnungsdatum.}
\newcommand{\unserebankverbindung} {Unsere Bankverbindung}
\newcommand{\textKontonummer} {Kontonummer:}
\newcommand{\textBank} {bei der}
\newcommand{\textBankleitzahl} {BLZ:}
\newcommand{\textBic} {BIC:}
\newcommand{\textIban} {IBAN:}
\newcommand{\unsereustid} {Unsere USt-Identifikationsnummer lautet}
\newcommand{\ihreustid} {Ihre USt-Identifikationsnummer:}
\newcommand{\steuerfreiEU} {Steuerfreie, innergemeinschaftliche Lieferung.}
\newcommand{\steuerfreiAUS} {Steuerfreie Lieferung ins außereuropäische Ausland.}

\newcommand{\textUstid} {UStId:}

% gutschrift (credit_note)
\newcommand{\gutschrift} {Gutschrift}
\newcommand{\fuerRechnung} {für Rechnung}
\newcommand{\gutschriftformel} {wir erlauben uns, Ihnen folgenden Positionen gutzuschreiben:}

% sammelrechnung (statement)
\newcommand{\sammelrechnung} {Statment}
\newcommand{\sammelrechnungsformel} {please note that following invoices are outstanding:}
\newcommand{\faellig} {Due}
\newcommand{\aktuell} {Current}
\newcommand{\asDreissig} {30}
\newcommand{\asSechzig} {60}
\newcommand{\asNeunzig} {90+}

% zahlungserinnerung (Mahnung)
\newcommand{\mahnung} {Zahlungserinnerung}
\newcommand{\mahnungsformel} {man kann seine Augen nicht überall haben - offensichtlich haben Sie übersehen, die folgenden Rechnungen zu begleichen:}
\newcommand{\beruecksichtigtBis} {Zahlungseingänge sind nur berücksichtigt bis zum}
\newcommand{\schonGezhalt} {Sollten Sie zwischenzeitlich bezahlt haben, betrachten Sie diese Zahlungserinnerung bitte als gegenstandslos.}

% zahlungserinnerung_invoice (Rechnung zur Mahnung)
\newcommand{\mahnungsrechnungsformel} {hiermit stellen wir Ihnen zu o.g.\ Zahlungserinerung folgende Posten in Rechnung:}
\newcommand{\posten} {Posten}
\newcommand{\betrag} {Betrag}
\newcommand{\bitteZahlenBis} {Bitte begleichen Sie diese Forderung bis zum}

% anfrage (request_quotion)
\newcommand{\anfrage} {Anfrage}
\newcommand{\anfrageformel} {bitte nennen Sie uns für folgende Artikel Preis und Liefertermin:}
\newcommand{\anfrageBenoetigtBis} {Wir benötigen die Lieferung bis zum}  %Danach wird das Datum eingefügt, falls das grammatisch nicht funktionieren sollte müssen wir eine ausnahme für die sprache definieren
\newcommand{\anfragedanke} {Im Voraus besten Dank für Ihre Bemühungen.}

% bestellung/auftrag (purchase_order)
\newcommand{\bestellung} {Bestellung}
\newcommand{\unsereBestellnummer} {Unsere Bestellnummer}
\newcommand{\bestellformel} {hiermit bestellen wir verbindlich folgende Positionen:}

% einkaufslieferschein (purchase_delivery_order)
\newcommand{\einkaufslieferschein} {Eingangslieferschein}
}{
	\ifstr{\lxlangcode}{DE}{
		\makeatletter
		\main@language{ngerman}
		\makeatother
		%%%%%%%%%%%%%%%%%%%%%%%%%%%%%%%%%%%%%%%%%%%%%%%
%standardphrasen und schnipsel in deutsch     %
%dient als vorlage für alle anderen sprachen  %
%%%%%%%%%%%%%%%%%%%%%%%%%%%%%%%%%%%%%%%%%%%%%%%


\newcommand{\anrede} {Sehr geehrte Damen und Herren,}
\newcommand{\anredefrau} {Sehr geehrte Frau}
\newcommand{\anredeherr} {Sehr geehrter Herr}


\newcommand{\nr} {Nr.}
\newcommand{\datum} {Datum}
\newcommand{\kundennummer} {Kunden-Nr.}
\newcommand{\ansprechpartner} {Ansprechpartner}
\newcommand{\bearbeiter} {Bearbeiter}
\newcommand{\gruesse} {Mit freundlichen Grüßen}
\newcommand{\vom} {vom}
\newcommand{\von} {von}
\newcommand{\seite} {Seite}
\newcommand{\uebertrag} {Übertrag}


\newcommand{\position} {Pos.}
\newcommand{\artikelnummer} {Art.-Nr.}
\newcommand{\bild} {Bild}
\newcommand{\keinbild} {kein Bild}
\newcommand{\menge} {Menge}
\newcommand{\bezeichnung} {Bezeichung}
\newcommand{\seriennummer}{Seriennummer}
\newcommand{\ean}{EAN}
\newcommand{\projektnummer}{Projektnummer}
\newcommand{\charge}{Charge}
\newcommand{\mhd}{MHD}
\newcommand{\einzelpreis} {E-Preis}
\newcommand{\gesamtpreis} {G-Preis}
\newcommand{\nettobetrag} {Nettobetrag}
\newcommand{\schlussbetrag} {Gesamtbetrag}

\newcommand{\weiteraufnaechsterseite} {weiter auf der nächsten Seite ...}

\newcommand{\zahlung} {Zahlungsbedingungen:}
\newcommand{\lieferung} {Lieferbedingungen:}
\newcommand{\textTelefon} {Tel.:}
\newcommand{\textFax} {Fax:}

% angebot (sales_quotion)
\newcommand{\angebot} {Angebot}
\newcommand{\angebotsformel} {gerne unterbreiten wir Ihnen folgendes Angebot:}
\newcommand{\angebotdanke} {Wir danken für Ihre Anfrage und hoffen, Ihnen hiermit ein interessantes Angebot gemacht zu haben.}
\newcommand{\angebotgueltig} {Das Angebot ist freibleibend gültig bis zum}		%Danach wird das Datum eingefügt, falls das grammatisch nicht funktionieren sollte müssen wir eine ausnahme für die sprache definieren
\newcommand{\angebotfragen} {Sollten Sie noch Fragen oder Änderungswünsche haben, können Sie uns gerne jederzeit unter den unten genannten Telefonnummern oder E-Mail-Adressen kontaktieren.}
\newcommand{\angebotagb} {Bei der Durchführung des Auftrags gelten unsere AGB, die wir Ihnen gerne zuschicken.}


% auftragbestätigung (sales_order)
\newcommand{\auftragsbestaetigung} {Auftragsbestätigung}
\newcommand{\auftragsnummer} {Auftrag-Nr.}
\newcommand{\ihreBestellnummer} {Ihre Bestellnummer}
\newcommand{\auftragsformel} {hiermit bestätigen wir Ihnen folgende Bestellpostionen:}
\newcommand{\lieferungErfolgtAm} {Die Lieferung erfolgt am} %Danach wird das Datum eingefügt, falls das grammatisch nicht funktionieren sollte müssen wir eine ausnahme für die sprache definieren
\newcommand{\auftragpruefen} {Bitte kontrollieren Sie alle Positionen auf Übereinstimmung mit Ihrer Bestellung! Teilen Sie Abweichungen innerhalb von 3 Tagen mit!}
\newcommand{\proformarechnung} {Proforma Rechnung}

% lieferschein (sales_delivery_order)
\newcommand{\lieferschein} {Lieferschein}

% rechnung (invoice)
\newcommand{\rechnung} {Rechnung}
\newcommand{\rechnungsdatum} {Rechnungsdatum}
\newcommand{\ihrebestellung} {Ihr Bestellung}
\newcommand{\lieferdatum} {Lieferdatum}
\newcommand{\rechnungsformel} {für unsere Leistungen erlauben wir uns, folgende Positionen in Rechnung zu stellen:}
\newcommand{\zwischensumme} {Zwischensumme}
\newcommand{\leistungsdatumGleichRechnungsdatum} {Das Leistungsdatum entspricht, soweit nicht anders angegeben, dem Rechnungsdatum.}
\newcommand{\unserebankverbindung} {Unsere Bankverbindung}
\newcommand{\textKontonummer} {Kontonummer:}
\newcommand{\textBank} {bei der}
\newcommand{\textBankleitzahl} {BLZ:}
\newcommand{\textBic} {BIC:}
\newcommand{\textIban} {IBAN:}
\newcommand{\unsereustid} {Unsere USt-Identifikationsnummer lautet}
\newcommand{\ihreustid} {Ihre USt-Identifikationsnummer:}
\newcommand{\steuerfreiEU} {Steuerfreie, innergemeinschaftliche Lieferung.}
\newcommand{\steuerfreiAUS} {Steuerfreie Lieferung ins außereuropäische Ausland.}

\newcommand{\textUstid} {UStId:}

% gutschrift (credit_note)
\newcommand{\gutschrift} {Gutschrift}
\newcommand{\fuerRechnung} {für Rechnung}
\newcommand{\gutschriftformel} {wir erlauben uns, Ihnen folgenden Positionen gutzuschreiben:}

% sammelrechnung (statement)
\newcommand{\sammelrechnung} {Sammelrechnung}
\newcommand{\sammelrechnungsformel} {bitte nehmen Sie zur Kenntnis, dass folgende Rechnungen unbeglichen sind:}
\newcommand{\faellig} {Fälligkeit}
\newcommand{\aktuell} {aktuell}
\newcommand{\asDreissig} {30}
\newcommand{\asSechzig} {60}
\newcommand{\asNeunzig} {90+}

% zahlungserinnerung (Mahnung)
\newcommand{\mahnung} {Zahlungserinnerung}
\newcommand{\mahnungsformel} {man kann seine Augen nicht überall haben - offensichtlich haben Sie übersehen, die folgenden Rechnungen zu begleichen:}
\newcommand{\beruecksichtigtBis} {Zahlungseingänge sind nur berücksichtigt bis zum}
\newcommand{\schonGezahlt} {Sollten Sie zwischenzeitlich bezahlt haben, betrachten Sie diese Zahlungserinnerung  bitte als gegenstandslos.}

% zahlungserinnerung_invoice (Rechnung zur Mahnung)
\newcommand{\mahnungsrechnungsformel} {hiermit stellen wir Ihnen zu o.g. \mahnung \ folgende Posten in Rechnung:}
\newcommand{\posten} {Posten}
\newcommand{\betrag} {Betrag}
\newcommand{\bitteZahlenBis} {Bitte begleichen Sie diese Forderung bis zum}

% anfrage (request_quotion)
\newcommand{\anfrage} {Anfrage}
\newcommand{\anfrageformel} {bitte nennen Sie uns für folgende Artikel Preis und Liefertermin:}
\newcommand{\anfrageBenoetigtBis} {Wir benötigen die Lieferung bis zum}  %Danach wird das Datum eingefügt, falls das grammatisch nicht funktionieren sollte müssen wir eine ausnahme für die sprache definieren
\newcommand{\anfragedanke} {Im Voraus besten Dank für Ihre Bemühungen.}

% bestellung/auftrag (purchase_order)
\newcommand{\bestellung} {Bestellung}
\newcommand{\unsereBestellnummer} {Unsere Bestellnummer}
\newcommand{\bestellformel} {hiermit bestellen wir verbindlich folgende Positionen:}

% einkaufslieferschein (purchase_delivery_order)
\newcommand{\einkaufslieferschein} {Eingangslieferschein}
}{%%%%%%%%%%%%%%%%%%%%%%%%%%%%%%%%%%%%%%%%%%%%%%%
%standardphrasen und schnipsel in deutsch     %
%dient als vorlage für alle anderen sprachen  %
%%%%%%%%%%%%%%%%%%%%%%%%%%%%%%%%%%%%%%%%%%%%%%%


\newcommand{\anrede} {Sehr geehrte Damen und Herren,}
\newcommand{\anredefrau} {Sehr geehrte Frau}
\newcommand{\anredeherr} {Sehr geehrter Herr}


\newcommand{\nr} {Nr.}
\newcommand{\datum} {Datum}
\newcommand{\kundennummer} {Kunden-Nr.}
\newcommand{\ansprechpartner} {Ansprechpartner}
\newcommand{\bearbeiter} {Bearbeiter}
\newcommand{\gruesse} {Mit freundlichen Grüßen}
\newcommand{\vom} {vom}
\newcommand{\von} {von}
\newcommand{\seite} {Seite}
\newcommand{\uebertrag} {Übertrag}


\newcommand{\position} {Pos.}
\newcommand{\artikelnummer} {Art.-Nr.}
\newcommand{\bild} {Bild}
\newcommand{\keinbild} {kein Bild}
\newcommand{\menge} {Menge}
\newcommand{\bezeichnung} {Bezeichung}
\newcommand{\seriennummer}{Seriennummer}
\newcommand{\ean}{EAN}
\newcommand{\projektnummer}{Projektnummer}
\newcommand{\charge}{Charge}
\newcommand{\mhd}{MHD}
\newcommand{\einzelpreis} {E-Preis}
\newcommand{\gesamtpreis} {G-Preis}
\newcommand{\nettobetrag} {Nettobetrag}
\newcommand{\schlussbetrag} {Gesamtbetrag}

\newcommand{\weiteraufnaechsterseite} {weiter auf der nächsten Seite ...}

\newcommand{\zahlung} {Zahlungsbedingungen:}
\newcommand{\lieferung} {Lieferbedingungen:}
\newcommand{\textTelefon} {Tel.:}
\newcommand{\textFax} {Fax:}

% angebot (sales_quotion)
\newcommand{\angebot} {Angebot}
\newcommand{\angebotsformel} {gerne unterbreiten wir Ihnen folgendes Angebot:}
\newcommand{\angebotdanke} {Wir danken für Ihre Anfrage und hoffen, Ihnen hiermit ein interessantes Angebot gemacht zu haben.}
\newcommand{\angebotgueltig} {Das Angebot ist freibleibend gültig bis zum}		%Danach wird das Datum eingefügt, falls das grammatisch nicht funktionieren sollte müssen wir eine ausnahme für die sprache definieren
\newcommand{\angebotfragen} {Sollten Sie noch Fragen oder Änderungswünsche haben, können Sie uns gerne jederzeit unter den unten genannten Telefonnummern oder E-Mail-Adressen kontaktieren.}
\newcommand{\angebotagb} {Bei der Durchführung des Auftrags gelten unsere AGB, die wir Ihnen gerne zuschicken.}


% auftragbestätigung (sales_order)
\newcommand{\auftragsbestaetigung} {Auftragsbestätigung}
\newcommand{\auftragsnummer} {Auftrag-Nr.}
\newcommand{\ihreBestellnummer} {Ihre Bestellnummer}
\newcommand{\auftragsformel} {hiermit bestätigen wir Ihnen folgende Bestellpostionen:}
\newcommand{\lieferungErfolgtAm} {Die Lieferung erfolgt am} %Danach wird das Datum eingefügt, falls das grammatisch nicht funktionieren sollte müssen wir eine ausnahme für die sprache definieren
\newcommand{\auftragpruefen} {Bitte kontrollieren Sie alle Positionen auf Übereinstimmung mit Ihrer Bestellung! Teilen Sie Abweichungen innerhalb von 3 Tagen mit!}
\newcommand{\proformarechnung} {Proforma Rechnung}

% lieferschein (sales_delivery_order)
\newcommand{\lieferschein} {Lieferschein}

% rechnung (invoice)
\newcommand{\rechnung} {Rechnung}
\newcommand{\rechnungsdatum} {Rechnungsdatum}
\newcommand{\ihrebestellung} {Ihr Bestellung}
\newcommand{\lieferdatum} {Lieferdatum}
\newcommand{\rechnungsformel} {für unsere Leistungen erlauben wir uns, folgende Positionen in Rechnung zu stellen:}
\newcommand{\zwischensumme} {Zwischensumme}
\newcommand{\leistungsdatumGleichRechnungsdatum} {Das Leistungsdatum entspricht, soweit nicht anders angegeben, dem Rechnungsdatum.}
\newcommand{\unserebankverbindung} {Unsere Bankverbindung}
\newcommand{\textKontonummer} {Kontonummer:}
\newcommand{\textBank} {bei der}
\newcommand{\textBankleitzahl} {BLZ:}
\newcommand{\textBic} {BIC:}
\newcommand{\textIban} {IBAN:}
\newcommand{\unsereustid} {Unsere USt-Identifikationsnummer lautet}
\newcommand{\ihreustid} {Ihre USt-Identifikationsnummer:}
\newcommand{\steuerfreiEU} {Steuerfreie, innergemeinschaftliche Lieferung.}
\newcommand{\steuerfreiAUS} {Steuerfreie Lieferung ins außereuropäische Ausland.}

\newcommand{\textUstid} {UStId:}

% gutschrift (credit_note)
\newcommand{\gutschrift} {Gutschrift}
\newcommand{\fuerRechnung} {für Rechnung}
\newcommand{\gutschriftformel} {wir erlauben uns, Ihnen folgenden Positionen gutzuschreiben:}

% sammelrechnung (statement)
\newcommand{\sammelrechnung} {Sammelrechnung}
\newcommand{\sammelrechnungsformel} {bitte nehmen Sie zur Kenntnis, dass folgende Rechnungen unbeglichen sind:}
\newcommand{\faellig} {Fälligkeit}
\newcommand{\aktuell} {aktuell}
\newcommand{\asDreissig} {30}
\newcommand{\asSechzig} {60}
\newcommand{\asNeunzig} {90+}

% zahlungserinnerung (Mahnung)
\newcommand{\mahnung} {Zahlungserinnerung}
\newcommand{\mahnungsformel} {man kann seine Augen nicht überall haben - offensichtlich haben Sie übersehen, die folgenden Rechnungen zu begleichen:}
\newcommand{\beruecksichtigtBis} {Zahlungseingänge sind nur berücksichtigt bis zum}
\newcommand{\schonGezahlt} {Sollten Sie zwischenzeitlich bezahlt haben, betrachten Sie diese Zahlungserinnerung  bitte als gegenstandslos.}

% zahlungserinnerung_invoice (Rechnung zur Mahnung)
\newcommand{\mahnungsrechnungsformel} {hiermit stellen wir Ihnen zu o.g. \mahnung \ folgende Posten in Rechnung:}
\newcommand{\posten} {Posten}
\newcommand{\betrag} {Betrag}
\newcommand{\bitteZahlenBis} {Bitte begleichen Sie diese Forderung bis zum}

% anfrage (request_quotion)
\newcommand{\anfrage} {Anfrage}
\newcommand{\anfrageformel} {bitte nennen Sie uns für folgende Artikel Preis und Liefertermin:}
\newcommand{\anfrageBenoetigtBis} {Wir benötigen die Lieferung bis zum}  %Danach wird das Datum eingefügt, falls das grammatisch nicht funktionieren sollte müssen wir eine ausnahme für die sprache definieren
\newcommand{\anfragedanke} {Im Voraus besten Dank für Ihre Bemühungen.}

% bestellung/auftrag (purchase_order)
\newcommand{\bestellung} {Bestellung}
\newcommand{\unsereBestellnummer} {Unsere Bestellnummer}
\newcommand{\bestellformel} {hiermit bestellen wir verbindlich folgende Positionen:}

% einkaufslieferschein (purchase_delivery_order)
\newcommand{\einkaufslieferschein} {Eingangslieferschein}
}
} % Ende EN


% Mandanten-/Firmenabhängigkeiten

% Pfad zu firmenspez. Angaben
% Hat man mehrere Mandanten muß man statt "Firma1" den Datenbanknamen seines
% Mandanten eingeben.

\ExplSyntaxOn
\int_set:Nn \l_kivi_tmp_int {1}
\bool_set_true:N \l_kivi_tmp_bool
\bool_while_do:Nn \l_kivi_tmp_bool {
	\file_if_exist:nTF {firma\int_use:N \l_kivi_tmp_int/ident.tex}
	{
	\str_if_in:NnTF \kivicompany {Firma\int_use:N \l_kivi_tmp_int}
		{
		\newcommand*{\identpath}{firma\int_use:N \l_kivi_tmpa_int}
		\bool_set_false:N \l_kivi_tmp_bool
		}
		{\int_incr:N \l_kivi_tmp_int}
	}
	{
	\bool_set_false:N \l_kivi_tmp_bool
	\newcommand*{\identpath}{firma}
	}
}

\ExplSyntaxOff


% Identität
\newcommand{\telefon} {++49 228 92 98 2012}
\newcommand{\fax} {}
\newcommand{\firma} {Richardson \& Büren GmbH}
\newcommand{\strasse} {Römerstr. 45 - 47}
\newcommand{\ort} {53111 Bonn}
\newcommand{\ustid} {DE292363254}
\newcommand{\finanzamt} {Finanzamt Bonn-Innenstadt}
\newcommand{\email} {information@kivitendo-premium.de}
\newcommand{\homepage} {http://www.kivitendo-premium.de}


\ExplSyntaxOn

%Setze Briefkopf-logo falls vorhanden
\setkomavar{fromlogo}{\includegraphics[width=.25\linewidth]{\identpath/briefkopf}}

% Währungen/Konten
\tl_new:N \g_kivi_currency_tl
\str_if_in:NnT \lxcurrency {USD} {\tl_gset:Nn \g_kivi_currency_tl {usd}}
\str_if_in:NnT \lxcurrency {CHF} {\tl_gset:Nn \g_kivi_currency_tl {chf}}
\str_if_in:NnT \lxcurrency {EUR} {\tl_gset:Nn \g_kivi_currency_tl {euro}}
\tl_if_empty:NT  \g_kivi_currency_tl {
	\tl_gset:Nn \g_kivi_currency_tl {default}
	\edef \currency {\tl_to_str:N \lxcurrency}
}

\input{\identpath/\g_kivi_currency_tl _account.tex}

\ExplSyntaxOff


% keine Absätze nach rechts einrücken
\setlength\parindent{0pt}



% Befehl f. normale Schriftart und -größe
\renewcommand*{\familydefault}{\sfdefault}
\KOMAoptions{fontsize=10pt}

% Einstellungen f. Kopf und Fuss
\pagestyle{kivitendo.letter}
%\setheadwidth[20mm]{page} % Kopfzeile nach rechts verschieben
%\setfootwidth[-39mm]{page} % Fusszeile verschieben

% Befehl f. laufende Kopfzeile:
% 1. Text f. Kunden- oder Lieferantennummer (oder leer, wenn diese nicht ausgegeben werden soll)
% 2. Kunden- oder Lieferantennummer (oder leer)
% 3. Belegname {oder leer}
% 4. Belegnummer {oder leer}
% 5. Belegdatum {oder leer}
% Beispiel: \ourhead{\kundennummer}{<%customernumber%>}{\angebot}{<%quonumber%>}{<%quodate%>}
\setkomafont{pagehead}{\scriptsize}
\newcommand{\ourhead}[5] {
\chead{
  \ifnum\thepage=1
  \else
      \ifstr{#1}{}{}{#1: #2 \hspace{0.7cm}}
      #3
      \ifstr{#4}{}{}{~\nr: #4}
      \ifstr{#5}{}{}{\vom ~ #5}
      \hspace{0.7cm} - \seite ~ \thepage/\letterlastpage  ~-%
  \fi
}
}

% Firmenfuss
\setkomafont{pagefoot}{\tiny}
\cfoot{
  {
     \begin{tabular}{p{5cm}p{4.5cm}lr}
        \firma                 & \email              & \textKontonummer & \kontonummer \\
        \strasse               & \homepage           & \textBank        & \bank \\
        \ort                   & \textUstid\ \ustid  & \textIban        & \iban \\
        \textTelefon~\telefon  & \finanzamt          & \textBic         & \bic \\
        \ifstr{\fax}{}{}{\textFax~\fax} & &\textBankleitzahl	& \bankleitzahl \\
     \end{tabular}
  }
}
